\documentclass[11pt]{article}
\usepackage[margin=1in]{geometry}
\usepackage{amsmath,amssymb,amsthm,enumitem,hyperref}
\hypersetup{colorlinks=true,linkcolor=blue,urlcolor=blue,citecolor=blue}
\newtheorem{theorem}{Theorem}
\newtheorem{lemma}{Lemma}
\newcommand{\Z}{\mathbb Z}
\newcommand{\Q}{\mathcal Q}
\title{Transpose duality for preorder multichain matrices\\\large MathDB 374957}
\author{Principia Math In Progress}
\date{September 18, 2026}
\begin{document}
\maketitle
\begin{abstract}
We prove the matrix Ehrhart--zeta duality conjecture of Athanasiadis
and Chapoton. Coordinate increments identify multichains with
lattice points in a sum of simplices attached to a bipartite graph
with repeated vertices. Exchanging the two vertex classes reverses
the preorder and exchanges the two parameters. Postnikov's
established lattice-point duality then gives the claimed transpose.
\end{abstract}

\section*{Task 1. Full solution}
Let $\tau$ be a preorder on a finite set $E$. Its preorder polytope is
\[
 \Q_\tau=\{x\in\mathbb R_{\ge0}^E:
          \sum_{e\in I}x_e\le |I|
          \text{ for every order ideal }I\text{ of }\tau\}.
\]
For $k,\ell\ge0$, let $\nabla_\tau(k,\ell)$ count sequences
\[
 a^{(1)}\le a^{(2)}\le\cdots\le a^{(k)},
 \qquad a^{(j)}\in\ell\Q_\tau\cap\Z^E,
\]
where comparison is coordinatewise. These are the $k$-multichains
of \cite[Section 4.2]{AC}. Let $\tau^*$ be the reversed preorder.

\begin{theorem}\label{thm:main}
For every finite preorder and all $k,\ell\ge0$,
\[
 \boxed{\nabla_\tau(k,\ell)=\nabla_{\tau^*}(\ell,k).}
\]
Thus Conjecture 4.2 of \cite{AC} holds.
\end{theorem}

\subsection*{The established bipartite-graph duality}
Let $G$ have disjoint vertex classes $A,B$, and let $N(S)$ be
the neighbors of a subset of either class in the opposite class.
For a set $J\subseteq B$, put
$\Delta_J=\operatorname{conv}(0,e_b:b\in J)$ and define
\[
 P_G=\sum_{a\in A}\Delta_{N(a)}\subseteq\mathbb R^B.
\]
Let $G^*$ be the graph with its two vertex classes exchanged.
We use Postnikov's lattice-point duality
\begin{equation}\label{eq:duality}
 \#(P_G\cap\Z^B)=\#(P_{G^*}\cap\Z^A),
\end{equation}
recorded explicitly in \cite[Lemmas 2.4--2.6]{D}.
We recall the characterization of its integer points and derive
\eqref{eq:duality} from the counting theorem used by the primary source.

First,
\begin{equation}\label{eq:hall}
 P_G\cap\Z^B=
 \left\{z\in\Z_{\ge0}^B:
          \sum_{b\in S}z_b\le |N(S)|\text{ for every }S\subseteq B\right\}.
\end{equation}
Necessity holds because a simplex summand contributes at most
one to the sum over $S$, and contributes zero unless its index
lies in $N(S)$. Conversely, replace each $b$ by $z_b$ copies.
Every subset of copies projecting to $S$ has at most
$\sum_{b\in S}z_b\le|N(S)|$ members. Hall's marriage theorem
therefore matches all copies to distinct vertices of $A$.
For each matched $a$, choose its matched vector $e_b$ from
$\Delta_{N(a)}$, and choose $0$ for unmatched $a$.
Their sum is $z$, proving \eqref{eq:hall}.

To see \eqref{eq:duality} from \cite[Theorem 2.1]{AC}, lift the
simplices to the coordinates $B\sqcup\{0\}$, using the sets
$I_a=N(a)\sqcup\{0\}$ with coefficient $1$, and adjoin the
universal set $I_0=B\sqcup\{0\}$ with coefficient $0$.
Projection forgetting coordinate $0$ is a bijection of integer
points, since the sum of all coordinates is fixed at $|A|$.
The draconian vectors in that theorem are
\[
 d_0+\sum_{a\in A}d_a=|B|,\qquad d_a\ge0,
 \qquad \sum_{a\in T}d_a\le |N(T)|\quad(T\subseteq A).
\]
Indeed, the inequalities for nonempty $T\subseteq A$ follow
after canceling the extra vertex $0$; those involving the
universal set follow from the total-sum condition. The inequality
for $T=A$ ensures that $d_0$ is uniquely determined and nonnegative.
At these coefficients, every binomial factor in Theorem 2.1
is $1$. The number of lattice points is therefore the number of
vectors $(d_a)_{a\in A}$ satisfying the displayed inequalities.
By \eqref{eq:hall} for $G^*$, this is exactly the right side of
\eqref{eq:duality}. This graph duality is an established input,
not a new claim of this solution.

\subsection*{The graph of repeated preorder elements}
Assume $k,\ell\ge1$. Form $G=G_\tau(k,\ell)$ with
\[
 A=E\times[\ell],\qquad B=E\times[k],
\]
and join $(e,s)\in A$ to $(f,t)\in B$ exactly when
$e\le_\tau f$. For $S\subseteq B$, let $\pi(S)\subseteq E$
be its projection. If $D(T)$ denotes the order ideal generated
by $T\subseteq E$, then
\[
 N(S)=D(\pi(S))\times[\ell].
\]
Consequently \eqref{eq:hall} reduces to
\begin{equation}\label{eq:increments}
 z_{f,t}\ge0,\qquad
 \sum_{f\in I}\sum_{t=1}^k z_{f,t}\le\ell|I|
 \quad(I\text{ an order ideal of }\tau).
\end{equation}
Here is the reduction in both directions. Taking $S=I\times[k]$
in \eqref{eq:hall} gives \eqref{eq:increments}, since $D(I)=I$.
Conversely, for arbitrary $S$, take $I=D(\pi(S))$.
Nonnegativity and $S\subseteq I\times[k]$ give
\[
 \sum_{(f,t)\in S}z_{f,t}
 \le\sum_{f\in I,t\in[k]}z_{f,t}
 \le\ell|I|=|N(S)|,
\]
as required.

An integer vector in \eqref{eq:increments} determines the chain
\begin{equation}\label{eq:chain}
 a_f^{(j)}=\sum_{t=1}^j z_{f,t}\qquad(1\le j\le k).
\end{equation}
The terminal vector belongs to $\ell\Q_\tau$ by
\eqref{eq:increments}. Every earlier vector does too, since
the polytope is closed under decreasing nonnegative coordinates.
Conversely, a multichain gives the nonnegative increments
$z_{f,t}=a_f^{(t)}-a_f^{(t-1)}$, where $a^{(0)}=0$.
These two constructions are inverse, so
\begin{equation}\label{eq:id}
 \nabla_\tau(k,\ell)=\#(P_{G_\tau(k,\ell)}\cap\Z^{E\times[k]}).
\end{equation}

\begin{proof}[Proof of Theorem~\ref{thm:main}]
Exchanging the two graph classes gives exactly
\[
 G_\tau(k,\ell)^*=G_{\tau^*}(\ell,k),
\]
because $e\le_\tau f$ is equivalent to $f\le_{\tau^*}e$.
Apply \eqref{eq:duality} and then \eqref{eq:id} to both graphs.
This proves the theorem for positive parameters. If $k=0$ there
is one empty multichain; if $\ell=0$, the lattice-point poset
contains only $0$ and again there is one multichain. These cases
give the same identity. The empty ground set, if allowed, also
gives the constant count $1$.
\end{proof}

\subsection*{Scope, attribution, and verification}
No antisymmetry assumption was used, so equivalent distinct
elements of the preorder are included. The argument proves every
matrix entry for arbitrary parameters. The conjecture and
polytope definitions are due to Athanasiadis and Chapoton.
Postnikov's counting theorem and Hall's theorem are established
inputs. The repeated-vertex graph and increment identification
are the reduction resolving the target.

All 26 pages of the current primary were read, along with the
relevant graph-duality statements in \cite{D}. Four exact
verification groups pass: direct chain enumeration, 9,725
transpose identities, independent graph and increment-vector
enumeration, and further graph examples. Bounded checks found
no prior full target claim or identified journal final. Review
is by the authoring agent only, without independent peer review
or proof-assistant certification.

\section*{Tasks 2--5. Proposed follow-up directions}
\begin{enumerate}[label=\textbf{Task \arabic*.},start=2,leftmargin=*]
\item Turn Postnikov's graph duality into an explicit algorithm
mapping individual multichains between the two sides.
\item Determine refinements that retain the ranks of the chain
entries and remain compatible with exchanging the graph classes.
\item Study versions with nonuniform numbers of copies of each
preorder element, allowing different chain capacities by coordinate.
\item Derive effective formulas for these matrices in families of
preorders with small width or bounded graph complexity.
\end{enumerate}
\begin{thebibliography}{9}
\raggedright
\bibitem{AC} C. A. Athanasiadis and F. Chapoton,
\emph{Polytopes and posets associated to preorders},
\href{https://arxiv.org/abs/2605.26916v1}{arXiv:2605.26916v1}, 2026.
\bibitem{D} Z. Dai, Q. Hou, Z. Liu, W. Thawinrak, and H. Wang,
\emph{Counting Lattice Points in Minkowski Sums of Cross Polytopes},
\href{https://arxiv.org/abs/2608.16037v2}{arXiv:2608.16037v2}, 2026.
\end{thebibliography}
\end{document}
